\chapter{The Landscape: A Knob With \texorpdfstring{$10^{500}$}{10\textasciicircum 500} Clicks}

\epi{Zero knobs on the front panel.\\ A warehouse out back with $10^{500}$ of them, and no map.}

\fig[1.0]{plate_landscape}{The landscape, drawn. Every valley is a possible vacuum, a universe with its own constants of nature. String theory's equations allow an astronomical number of them, and nothing yet says which valley is ours.}


\section{A theory with no knobs at all}

After eight chapters of counting knobs in \LCDM's kitchen, you might expect string theory to be the opposite case, and in one sense it is. Its admirers point out, rightly, that string theory has \emph{no} adjustable dimensionless parameters in its fundamental formulation. There's a string length, which sets the units, and that's it. The coupling strength isn't a knob; it's the value of a field. Every number in physics, in principle, is supposed to come out of the theory.

That's an extraordinary ambition, and it's worth pausing to admire it. The Standard Model of particle physics has around nineteen free parameters, and more once you include neutrino masses: particle masses, mixing angles and coupling strengths, all measured and none explained. \LCDM\ adds its six. A theory that derived all of them from nothing but geometry would be the greatest achievement in the history of science.

The catch is in the phrase ``in principle.''

\section{Where the knobs went}

Superstring theory is consistent in ten spacetime dimensions. We see four. The other six must be curled up, \emph{compactified}, into a tiny space, and the physics we see in four dimensions depends on the shape and size of that space. The favoured candidates are Calabi--Yau manifolds, six-dimensional spaces with special properties that preserve some supersymmetry.

\fig{garden_hose}{How a dimension can hide. Up close, a garden hose is a two-dimensional surface; from far away it looks like a one-dimensional line. String theory curls up six extra dimensions this way, far too small to see.}



How many Calabi--Yau shapes are there? Nobody knows if the number is finite. One classical construction, as hypersurfaces in toric varieties, uses four-dimensional reflexive polytopes; Kreuzer and Skarke enumerated those in 2000 and found 473,800,776. Each gives at least one Calabi--Yau space, and many give more.

\fig{counts_ladder}{How big the numbers are, measured by how many digits it takes to write them: about nine for the count of four-dimensional reflexive polytopes, 500 for the famous estimate of flux vacua, and 272,000 for one F-theory geometry.}



Each shape also comes with continuous deformations, called \emph{moduli}: sizes and shape parameters that can vary smoothly without changing the topology. Their counts are given by the Hodge numbers $h^{1,1}$ and $h^{2,1}$. The quintic threefold, the simplest textbook example, has $h^{2,1}=101$ complex-structure moduli. In four dimensions every modulus shows up as a massless scalar field, which would mediate a long-range force that experiments don't see. So the moduli have to be fixed, or \emph{stabilised}.

That's where the knobs come from. The theory doesn't have them in its equations. It has them in its \emph{solutions}.

\section{Fluxes: how to build a discretuum}

The standard way to stabilise moduli is with fluxes: generalised magnetic fields threading the cycles of the compact space. Like magnetic flux through a superconducting ring, these fluxes are quantised. Each cycle carries an integer $n_i$.

Bousso and Polchinski showed in 2000 how this solves, or rather dissolves, the cosmological constant problem. Suppose there are $J$ independent cycles with flux quanta $n_1,\dots,n_J$ and charges $q_i$. The effective cosmological constant is roughly
\begin{equation}
\Lambda=\Lambda_{\rm bare}+\frac12\sum_{i=1}^{J}n_i^2q_i^2,
\end{equation}
with $\Lambda_{\rm bare}$ a large negative number. Each choice of integers $(n_1,\dots,n_J)$ is a different vacuum, a different universe with its own physics.

\fig{flux_lattice}{The Bousso--Polchinski picture in two dimensions. Each lattice point is a choice of flux integers. The cosmological constant is near zero only for points in a thin shell (red). In two dimensions there are a handful; in hundreds of dimensions there are astronomically many.}


\begin{mathbox}{Counting vacua in $J$ dimensions}
Think of the flux vector $\mathbf{n}=(n_1q_1,\dots,n_Jq_J)$ as a point in a $J$-dimensional lattice with cells of size about $q$. A vacuum has $\Lambda$ near zero when $|\mathbf n|^2\approx2|\Lambda_{\rm bare}|$, a sphere of radius $r=\sqrt{2|\Lambda_{\rm bare}|}$. The number of lattice points inside a ball of that radius is about
\[
\mathcal{N}\approx\frac{\pi^{J/2}}{\Gamma(J/2+1)}\left(\frac{r}{q}\right)^{J}\approx\frac{1}{\sqrt{\pi J}}\left(\frac{2\pi e\,r^2}{J\,q^2}\right)^{J/2},
\]
using Stirling's formula for the Gamma function. The count is astronomical as soon as $r/q$ exceeds about $\sqrt{J/2\pi e}$. With $J=500$ cycles, fluxes up to about $\pm10$ in units of their charges give $\mathcal N\approx10^{132}$, and fluxes up to about $\pm54$ give $\mathcal N\approx10^{500}$. The values of $\Lambda$ for these vacua are spread over a range of order the Planck scale, so the typical spacing between neighbouring values is about $1/\mathcal N$ in Planck units, far below the observed value of about $10^{-120}$. Somewhere in the lattice there are vacua with a cosmological constant as small as ours.
\end{mathbox}

\fig{vacua_count}{How many flux vacua? The ball count as the number of cycles $J$ grows, for three ranges of flux. Above the gold line there are enough vacua for some to have a cosmological constant as small as ours; with a few hundred cycles and fluxes of a few tens, the count reaches $10^{500}$.}

\fig{powers_of_ten_zoom}{What one part in $10^{120}$ means. Zoom in by ten on a ruler, then by ten again, and again. Three zooms are shown; finding one setting in $10^{120}$ takes 120 of them.}



That's the origin of the famous number. Later work refined it. Douglas and collaborators developed statistical methods for counting flux vacua, in 2003 Kachru, Kallosh, Linde and Trivedi (KKLT) proposed a construction of metastable vacua with a small positive cosmological constant, and the estimate of ``$10^{500}$'' entered folklore. In 2015 Taylor and Wang found a single F-theory geometry with of order $10^{272{,}000}$ flux vacua, which makes $10^{500}$ look modest.

\section{The landscape and the anthropic principle}

Leonard Susskind named this the \emph{landscape}, in 2003, and argued that we should embrace it. If there are $10^{500}$ vacua, and if cosmic inflation populates many of them in a vast multiverse, then the cosmological constant in our region needn't be explained by any equation. It's small because we live in one of the rare places where it's small enough for galaxies to form and observers to exist.

This has a respectable ancestor. In 1987 Steven Weinberg argued that if the vacuum energy were more than a few hundred times today's matter density, gravitational clumping would stall before galaxies could form, and that observers should therefore expect to find $\Lambda$ not too far below that bound. A decade later $\Lambda$ was measured, within an order of magnitude or so of the region his argument favoured. That's the best anthropic argument in physics, and it deserves the credit.

\fig{weinberg_window}{Weinberg's window. On a scale running up to the naive Planck-scale estimate, the range where galaxies can form (green) is a sliver near zero, and the measured value sits inside it (right panel, zoomed).}



But notice what happened to the knob. In \LCDM, $\Lambda$ is a dial you read off the sky. In the landscape, $\Lambda$ is one of $10^{500}$ settings of a dial nobody can turn, selected by the fact that someone is around to read it.

\section{A knob measured in bits}

Chapter~1 gave a way to put every knob on one scale: the information needed to pin it down. Let's do that for the landscape.

Choosing one vacuum out of $10^{500}$ requires
\begin{equation}
I_{\rm landscape}=\log_2\left(10^{500}\right)=500\log_210\approx1661\ \text{bits}.
\end{equation}
For the Taylor--Wang geometry it's $272{,}000\log_210\approx903{,}564$ bits.

Compare:
\begin{center}
\begin{tabular}{lr}
\toprule
knob & bits\\
\midrule
my galaxy law, $\kappa$ to 5\% over a 2-dex prior & 6.6\\
base \LCDM, 6 knobs at Planck precision & $\sim50$\\
the landscape, $10^{500}$ vacua & 1661\\
one F-theory geometry, $10^{272{,}000}$ & 903,564\\
\bottomrule
\end{tabular}
\end{center}

\fig{bits_ladder}{Knobs in bits, on a logarithmic scale. My law's one knob, \LCDM's six, everything physics has ever measured, and the landscape.}


\begin{mathbox}{A Fermi estimate: can we even find our vacuum?}
To \emph{predict} anything from the landscape, you'd first locate our vacuum by matching its low-energy physics to what we measure. How much information do our measurements contain?

The Standard Model and cosmology have roughly 30 measured constants. A few are known extremely well: the fine-structure constant to about a tenth of a part per billion, which is about $\log_2(10^{10})\approx33$ bits. Most are known to between a tenth of a per cent and ten per cent, which is 3 to 10 bits. Be generous and call it 12 bits each on average:
\[
I_{\rm measured}\approx30\times12\approx360\ \text{bits}.
\]
If vacua were spread uniformly and independently over those constants, our measurements would narrow $10^{500}$ down to about
\[
2^{1661-360}=2^{1301}\approx10^{392}\ \text{vacua}
\]
that agree with everything we have ever measured. This is cartoonish. Vacua aren't uniform, constants aren't independent, and most vacua don't look anything like the Standard Model. But the order of magnitude makes the point: the landscape's selection problem is bigger than all the information physics has ever extracted from nature, by hundreds of orders of magnitude.
\end{mathbox}

\section{The case for the defence}

It would be cheap to stop there, and it would also be unfair, so here's the defence.

\textbf{It's not a fit.} The landscape isn't a set of knobs someone tunes to match data. It's a claim about the solution space of a fixed theory. A knob in a \LCDM\ recipe is freedom chosen by a modeller; a vacuum in the landscape is a solution the theory forces on you. You can't simply decline to have it.

\textbf{It produced real physics.} String theory gave the first microscopic counting of black hole entropy (Strominger and Vafa, 1996), the AdS/CFT correspondence that has transformed how physicists think about quantum gravity and strongly coupled systems, and a great deal of mathematics. None of that depends on the landscape being small.

\textbf{People are trying to cut it down.} The \emph{swampland} programme, begun by Vafa in 2005, asks which effective theories \emph{can't} come from string theory. One conjecture, from Obied, Ooguri, Spodyneiko and Vafa in 2018, proposed that stable de Sitter vacua, universes with a true positive cosmological constant, may not exist in string theory at all. If true, dark energy would have to be evolving. You can imagine the interest in that community when DESI's evolving-dark-energy hint arrived.

\textbf{The critics are also physicists.} Lee Smolin's \emph{The Trouble with Physics} and Peter Woit's \emph{Not Even Wrong} (both 2006) argued that a theory compatible with almost any low-energy physics has stopped making testable predictions. In 2014 George Ellis and Joe Silk wrote in \emph{Nature} that some theorists were proposing to weaken the requirement of testability, and called on physicists to defend it. The debate is ongoing and conducted, for the most part, in good faith.

\section{The punch}

Here's the honest jab, in the spirit of this book. A theory with $10^{500}$ vacua has, in effect, one knob with $10^{500}$ clicks. It's the most knobby theory ever proposed, wearing the costume of a theory with no knobs at all. It moves the freedom from the control panel to the warehouse, and the warehouse has no index.

And the anthropic escape converts that freedom into an argument that we shouldn't expect explanations for the things we most want explained, like the value of $\Lambda$. It's possible that this is how the universe is. If so, it's a sad day for knob accounting, because it means some knobs can never be counted.

\begin{knobbox}{Knob accounting for the landscape}
\begin{itemize}
\item Front-panel knobs: zero dimensionless parameters in the fundamental theory.
\item Warehouse knobs: one choice of vacuum among $\sim10^{500}$ (or $10^{272{,}000}$ in one geometry).
\item Information needed to select: $\sim1661$ bits, against $\sim360$ bits in all measured constants.
\item Unique low-energy predictions: none established.
\end{itemize}
\end{knobbox}

\section{The question}

We've now seen three control panels.

\LCDM\ at the CMB: six knobs, superbly set, with an unexplained cosmological constant. \LCDM\ in galaxies: twelve knobs to imitate a formula closely, about 1066 to imitate it and its neighbours all at once, and five hand terms to make it notice dark energy. The landscape: a single knob with $10^{500}$ clicks and no map.

So here's the question this book has been building toward. What about \emph{one} knob? One number, set once, for every disc galaxy in the universe, with its value tied to the dark energy density measured by the expansion of the universe itself?

What does one knob get right? Where does it fail? And is there a test it can't wriggle out of?

That's Part~IV.
